\begin{proofbox}
\textbf{方法一：几何投影法——看清四项从哪里来}

先以 $B$ 为局部原点，取 $u,v>0,\ 0^\circ<\theta<90^\circ$。
令 $H$ 为水平段旋转后的终点，旋转前后的两段始终垂直。
\begin{center}\begin{tikzpicture}[x=1cm,y=1cm,>=Stealth,font=\small]
\path[use as bounding box] (-6.6,-0.6) rectangle (6.6,5.2);

\node at (0,4.8) {先转两段路，再把横、纵变化相加（图取 $u=3,v=2,\theta=30^\circ$）};
\begin{scope}[shift={(-4.8,0.3)}]
\draw[->,gray] (-0.4,0)--(3.7,0) node[right] {$u$方向};
\draw[->,gray] (0,-0.3)--(0,3.2);
\draw[blue!75!black,very thick,->] (0,0)--(3,0) node[midway,below] {$u$};
\draw[orange!85!black,very thick,->] (3,0)--(3,2) node[midway,right] {$v$};
\draw[gray,dashed] (0,0)--(3,2);
\node[below left] at (0,0) {$B$}; \node[above] at (3,2) {$P$};
\draw (2.8,0)--(2.8,0.2)--(3,0.2);
\end{scope}
\begin{scope}[shift={(1.5,0.3)}]
\coordinate (H) at ({3*cos(30)},{3*sin(30)});
\coordinate (Q) at ({3*cos(30)-2*sin(30)},{3*sin(30)+2*cos(30)});
\draw[->,gray] (-0.5,0)--(4,0) node[right] {$x$};
\draw[->,gray] (0,-0.3)--(0,3.9) node[above] {$y$};
\draw[blue!75!black,very thick,->] (0,0)--(H);
\draw[orange!85!black,very thick,->] (H)--(Q);
\draw[teal!75!black,dashed,thick] (0,0)--(Q);
\draw[blue!75!black,densely dotted] (H)--({3*cos(30)},0);
\node[blue!75!black,below] at (1.3,0) {$u\cos\theta$};
\node[blue!75!black,right] at (2.6,0.7) {$u\sin\theta$};
\draw[orange!85!black,densely dotted] (Q)--({3*cos(30)},{3*sin(30)+2*cos(30)})--(H);
\node[orange!85!black,above] at (2.1,3.65) {$-v\sin\theta$};
\node[orange!85!black,right] at (2.6,2.5) {$v\cos\theta$};
\draw[->] (0.65,0) arc(0:30:0.65); \node at (0.95,0.35) {$\theta$};
\fill (H) circle(1.5pt) node[below right] {$H$};
\fill[teal!75!black] (Q) circle(2pt) node[above left] {$Q$};
\node[below left] at (0,0) {$B$};
\draw ($(H)+(-0.173,-0.1)$)--($(H)+(-0.273,0.073)$)--($(H)+(-0.1,0.173)$);
\end{scope}

\end{tikzpicture}\end{center}

蓝色段长度为 $u$，与水平方向夹角为 $\theta$，故横、纵变化为
$u\cos\theta,u\sin\theta$。橙色段长度为 $v$，与竖直方向夹角为 $\theta$，
故向左走 $v\sin\theta$、向上走 $v\cos\theta$。于是
\[
s=u\cos\theta-v\sin\theta,\qquad
t=u\sin\theta+v\cos\theta.
\]
\textbf{如何覆盖所有角与坐标？}
由任意角三角函数的坐标定义，长度为 $1$、方向角为 $\theta$ 的有向线段坐标为
$(\cos\theta,\sin\theta)$；与其逆时针垂直的单位线段方向角为 $\theta+90^\circ$，
坐标为 $(-\sin\theta,\cos\theta)$。两段分别乘以带符号的 $u,v$ 后相加，仍得上式。
因此不需要把超出锐角的转动猜成正长度投影。最后加回 $(a,b)$，即得原结论。
\end{proofbox}
\clearpage
\begin{proofbox}
\textbf{方法二：拆箭头法——旋转两个单位方向}
\begin{center}\begin{tikzpicture}[x=1cm,y=1cm,>=Stealth,font=\small]
\path[use as bounding box] (-6.6,-0.6) rectangle (6.6,5.1);

\node at (0,4.7) {两个单位方向一起转 $\theta$，仍然垂直且长度为 $1$};
\begin{scope}[shift={(-3.7,1.3)},scale=2]
\draw[->,gray] (-1.3,0)--(1.5,0) node[right] {$x$};
\draw[->,gray] (0,-0.5)--(0,1.4) node[above] {$y$};
\draw[gray,dashed] (1,0) arc(0:180:1);
\draw[blue!75!black,->,very thick] (0,0)--(1,0) node[below] {$(1,0)$};
\draw[orange!85!black,->,very thick] (0,0)--(0,1) node[above left] {$(0,1)$};
\draw[teal!75!black,->,very thick] (0,0)--(30:1);
\draw[purple,->,very thick] (0,0)--(120:1);
\draw[->] (0.4,0) arc(0:30:0.4);
\node at (20:0.62) {$\theta$};
\end{scope}
\node[anchor=west,teal!75!black] at (0.4,3.3) {横向单位箭头 $\to(\cos\theta,\sin\theta)$};
\node[anchor=west,purple] at (0.4,2.3) {纵向单位箭头 $\to(-\sin\theta,\cos\theta)$};
\node[anchor=west,align=left] at (0.4,0.8) {先分别乘 $u,v$\\再首尾相接求和};

\end{tikzpicture}\end{center}

设旋转前的两个单位方向为 $\mathbf e_1=(1,0),\ \mathbf e_2=(0,1)$。
由单位圆定义，它们转后分别为
\[
\mathbf f_1=(\cos\theta,\sin\theta),\quad
\mathbf f_2=(\cos(\theta+90^\circ),\sin(\theta+90^\circ))
=(-\sin\theta,\cos\theta).
\]
原来 $\overrightarrow{BP}=u\mathbf e_1+v\mathbf e_2$。
\textbf{刚性旋转保持平行、长度和首尾相接关系}，所以原来的两段移动转后分别为
$u\mathbf f_1$ 和 $v\mathbf f_2$，包括 $u,v$ 为负或零的情形。故
\[
\begin{aligned}
\overrightarrow{BQ}&=u\mathbf f_1+v\mathbf f_2\\
&=(u\cos\theta-v\sin\theta,\ u\sin\theta+v\cos\theta).
\end{aligned}
\]
再加上中心坐标，得到 $Q$。此法用两支单位箭头概括了方法一的两段投影。
\end{proofbox}
\clearpage
\begin{proofbox}
\textbf{方法三：向量坐标法——两个垂直方向给出两条等式}

记 $(u,v)=\overrightarrow{BP},\ (s,t)=\overrightarrow{BQ}$，先设 $BP=BQ=r>0$。
\begin{center}\begin{tikzpicture}[x=1cm,y=1cm,>=Stealth,font=\small]
\path[use as bounding box] (-6.6,-0.6) rectangle (6.6,5.3);

\node at (0,4.9) {沿 $\overrightarrow{BP}$ 与其逆时针垂直方向分解 $\overrightarrow{BQ}$};
\begin{scope}[shift={(-3.8,0.6)},scale=0.85]
\coordinate (P) at (4,0); \coordinate (Q) at (50:4); \coordinate (F) at ({4*cos(50)},0);
\draw[->,gray] (-0.6,0)--(5.5,0);
\draw[->,gray] (0,-0.3)--(0,4.2) node[above left] {$(-v,u)$方向};
\draw[blue!75!black,very thick,->] (0,0)--(P) node[below] {$P$};
\draw[teal!75!black,very thick,->] (0,0)--(Q) node[above] {$Q$};
\draw[orange!85!black,dashed,thick] (Q)--(F);
\node[below] at (1.28,0) {$r\cos\theta$};
\node[right] at (2.57,1.6) {$r\sin\theta$};
\draw ($(F)+(-0.2,0)$)--($(F)+(-0.2,0.2)$)--($(F)+(0,0.2)$);
\draw[->] (0.7,0) arc(0:50:0.7); \node at (25:1.05) {$\theta$};
\node[below left] at (0,0) {$B$};
\node[blue!75!black] at (4.4,-0.65) {$(u,v)$方向};
\end{scope}
\node[anchor=west,align=left] at (1.8,3.5) {图中 $BP=BQ=r$\\两条轴为局部方向\\不一定是原 $x,y$ 轴};
\node[anchor=west,align=left] at (1.8,1.4) {沿原向量：$us+vt$\\沿垂直方向：$ut-vs$};

\end{tikzpicture}\end{center}

以 $\mathbf p=(u,v)$ 为第一个方向，以 $\mathbf n=(-v,u)$ 为第二个方向。
两者垂直且长度都是 $r$。图中沿两个方向的有向投影分别为 $r\cos\theta,r\sin\theta$，故由数量积得
\[
\mathbf p\cdot(s,t)=us+vt=r^2\cos\theta,\qquad
\mathbf n\cdot(s,t)=ut-vs=r^2\sin\theta.
\]
第二式也可由 $\mathbf n$ 的方向比 $\mathbf p$ 多 $90^\circ$ 看出：
$\cos(\theta-90^\circ)=\sin\theta$，因此适用于任意角。

第一式乘 $u$，减去第二式乘 $v$，得
$(u^2+v^2)s=r^2(u\cos\theta-v\sin\theta)$；
第一式乘 $v$，加上第二式乘 $u$，得
$(u^2+v^2)t=r^2(v\cos\theta+u\sin\theta)$。
因 $u^2+v^2=r^2>0$，除去 $r^2$ 即得 $s,t$，且解唯一。
$r=0$ 时 $P=B=Q$，公式直接成立。最后加回中心即可。

\textbf{反向使用：} 两条数量积等式除以 $r^2$，就是反求 $\cos\theta,\sin\theta$ 的公式。
\end{proofbox}
\clearpage
\begin{proofbox}
\textbf{方法四：三角函数法——方向角相加}
\begin{center}\begin{tikzpicture}[x=1cm,y=1cm,>=Stealth,font=\small]
\path[use as bounding box] (-6.6,-0.6) rectangle (6.6,5.3);

\node at (0,4.9) {半径不变，方向角由 $\alpha$ 变为 $\alpha+\theta$};
\begin{scope}[shift={(-3.6,0.4)},scale=0.95]
\draw[->,gray] (-0.4,0)--(4.8,0) node[right] {$x$};
\draw[->,gray] (0,-0.3)--(0,4.5) node[above] {$y$};
\draw[gray,dashed] (3.5,0) arc(0:110:3.5);
\draw[blue!75!black,very thick,->] (0,0)--(20:3.5) node[right] {$P$};
\draw[teal!75!black,very thick,->] (0,0)--(65:3.5) node[above right] {$Q$};
\draw[blue!75!black,->] (1,0) arc(0:20:1); \node at (10:1.4) {$\alpha$};
\draw[orange!85!black,->] (20:1.9) arc(20:65:1.9); \node at (43:2.3) {$\theta$};
\draw[teal!75!black,->] (0.65,0) arc(0:65:0.65);
\node[below left] at (0,0) {$B$};
\node[blue!75!black] at (2.5,0.6) {$r$};
\node[teal!75!black] at (0.7,2.5) {$r$};
\end{scope}
\node[anchor=west,align=left] at (1.6,3.2) {$P$：$(r\cos\alpha,r\sin\alpha)$\\[6pt]$Q$：$(r\cos(\alpha+\theta),$\\\hspace{1.1cm}$r\sin(\alpha+\theta))$};
\node[anchor=west] at (1.6,0.7) {这里均为相对 $B$ 的坐标};

\end{tikzpicture}\end{center}

若 $r=BP>0$，设 $\overrightarrow{BP}$ 的方向角为 $\alpha$，则
\[
u=r\cos\alpha,\quad v=r\sin\alpha,\qquad
s=r\cos(\alpha+\theta),\quad t=r\sin(\alpha+\theta).
\]
使用两角和公式展开：
\[
\begin{aligned}
s&=r\cos\alpha\cos\theta-r\sin\alpha\sin\theta
=u\cos\theta-v\sin\theta,\\
t&=r\sin\alpha\cos\theta+r\cos\alpha\sin\theta
=v\cos\theta+u\sin\theta.
\end{aligned}
\]
角 $\alpha,\theta$ 可以任意，故展开不受象限限制。若 $r=0$，两组相对坐标都为零，
公式仍成立；最后加回 $(a,b)$ 即完成证明。顺时针只需令 $\theta<0$。

\textbf{补充验算（不是另一种独立证明）：}
将所得 $s,t$ 代入并展开，有
\[
s^2+t^2=u^2+v^2,\quad
us+vt=(u^2+v^2)\cos\theta,\quad
ut-vs=(u^2+v^2)\sin\theta.
\]
三式依次核对长度、余弦和有向角的正弦。仅检验长度不能排除转向写反。
\end{proofbox}
